\section{A sharper adverse labeling and a test of its surviving vertex}
\label{sec:rootmark}

The best adverse labeling in the second search gives
\[
 \sup_{\text{rules}}\inf_{\text{labelings}}
 \frac{\mathcal R}{\mathcal K}
 \ \le\ \frac{1018318}{10^7}=10.18318\%.
\]
Here a rule assigns nonnegative rates of total 64 to the labeled residues. These rates may depend on the labeling, remain fixed across \(0<\varepsilon<1/52\), and meet the limiting marked-tile product certificate, including its source-window prefactor if used. Closures of this class satisfy the same upper bound. Paired with the constructive \(10.1190251313\%\) lower result for the prefactor version, the interval before the final census is about \(0.06415487\) \emph{percentage points}. This bound applies to the quotient experiment. The original selector could admit laws that pass by a different argument.

\subsection{Two consequences of using one profile across the source scales}

For an occupied coset of width \(w\), let \(M\) denote the sum of the rates in it and \(E\) the sum in its preferred child. Thus \(\mu=M/w\) and \(\nu=E/w\). The familiar sufficient condition is
\(10e^{-\mu}\nu^{3r/(6r+1)}>1\).
The refined union bound for the literal number of trial tiles allows its \((6r+1)\)-st power to exceed \((h+2)/(h+3)\). In the original scales \(h=2+8\varepsilon\) and \(r=\lceil32/\varepsilon+64\rceil+1\). As \(\varepsilon\) tends to zero, \(r\) tends to infinity, \(3r/(6r+1)\) tends to \(1/2\), and this positive prefactor raised to \(1/(6r+1)\) tends to one. Testing a fixed vector for every \(\varepsilon\) therefore gives the necessary closed inequality
\[
 E\ \ge\ \frac{w}{100}\exp(2M/w).
 \tag{R1}
\]
This is a necessity for either stated certificate across the scales. It does not assert the inequality for a vector allowed to change with \(\varepsilon\) in the refined test.
\href{https://github.com/nousentllc/riemann-hypothesis/blob/40e56bab026552e089d7afcfeff08d2defb9c005/src/Development/AutomorphicArithmetic/ThickFrames/MarkedLow.lean#L20-L31}{Original parameter definitions}.
\href{https://github.com/nousentllc/riemann-hypothesis/blob/40e56bab026552e089d7afcfeff08d2defb9c005/src/Formalism/SharpRestriction/GrowthBounds/UnramifiedQuadraticSetting.lean#L237-L251}{Trial-block multiplicity}.

At the root \(M=64\), so \(E\ge0.64e^2>4.72\). In a homogeneous positive leaf with rate \(a_t\), the finer cosets give \(e^{2a_t}\le50a_t\). The ratio \(e^{2a}/a\) increases for \(a>1/2\). The fortieth-term Taylor partial sum at \(2(23921/10000)\), even with its positive remainder removed, exceeds \(50(23921/10000)\). It follows that \(a_t\le C=23921/10000\) for every leaf, including zero-rate leaves. In particular at least two preferred root leaves must carry positive mass. This observation removes an entire class of possible supports when the root itself has a missing increment. A support maximizing a relaxation that ignores root holes may still be included with fewer preferred leaves; this only enlarges that relaxation.

\subsection{Detecting occupancy by the mass of a specific leaf}

The adverse word is \(2549776587809433532\), with bit \(s-1+c\) specifying the preferred child of residue \(c\bmod s\). Its root prefers the even numbers; inside the odd coset the preferred child consists of the residues \(3\bmod4\). Consider the width-32 odd coset. If it is empty, its preferred mass is zero. In an occupied coset the exponential in (R1) has a positive tangent intercept for \(u<1/2\), so applying its tangent at such \(u\) to an empty coset would be unjustified. A leaf provides a continuous lower bound on the occupancy indicator. If \(t\) is in this coset, then
\[
 \mathbf1_{\{M>0\}}\ \ge\ a_t/C.
\]
Fix \(u=3/8\), \(t=9\), and put \(k=\lfloor 10^9 S_{39}(2u)/100\rfloor/10^9\), where \(S_{39}\) sums through power 39 in the exponential series. Then every admissible rate vector satisfies
\[
 L(a)= E-2kM-32k(1-2u)\,a_9/C\ \ge\ 0.
 \tag{R2}
\]
Indeed, on a nonempty node, bounding the exponential below by its tangent gives \(E\ge \exp(2u)[2M+32(1-2u)]/100\). Both coefficients here are nonnegative, so one may round \(k\) down and then replace one by \(a_9/C\). On an empty node all terms in (R2) vanish. The choice of \(u\) is explained by the obstructing vertex below. A single price for the preferred \emph{root} mass had little force on that vertex. Its violation is one level lower.

For any nonnegative \(\ell\), the priced objective \(P-zQ+\ell L(a)\) upper bounds \(P-zQ\) on every feasible profile. Using \(\ell=803133/10^6\), we round each coefficient of \(\ell L\) upward to units of \(10^{-7}\). All masses are nonnegative, so rounding preserves the upper direction. Enumerate the full binary-tree support family, with the root elimination above, and relax each support to a closed capped simplex. The maximum of a fixed affine function on one of these simplices places \(C\) at 26 sites, \(1.8054=64-26C\) at one further site, and zero at the others. In particular the maximization can overcount a profile by allowing a nominally occupied leaf to receive zero mass. This freedom is needed for an upper bound without an artificial minimum positive rate.

\begin{center}
\begin{tabular}{@{}lrr@{}}
\toprule
Test at the final word & Numerator of \(z\) & Scaled maximum\\
\midrule
Target, unpriced  & 1018318 & \(54\,250\,880\,988\)\\
Target, one priced node & 1018318 & \(-22\,184\,178\)\\
Previous grid point, same price & 1018317 & \(8\,175\,060\)\\
\bottomrule
\end{tabular}
\end{center}
The score units are \(10^{-11}\), from combining the \(10^{-7}\) coefficient grid and \(10^{-4}\) mass units. A negative maximum certifies the bound for the continuum of rate vectors on each support. The positive maximum one tick earlier identifies the resolution of this particular price and relaxation. It gives no matching lower construction. Three comparison words from the same search yielded bounds \(10.18996\%\), \(10.20123\%\), and \(10.56603\%\) under their recorded choices of prices. The last of these used zero prices. A weaker closed-box optimum at one labeling was therefore a poor guide to which labeled constraint would produce the sharpest result after pricing. An independent integer census gave the same two signed maxima. It used 3,263,443 admissible masks for each 32-leaf half, and found 39,058 distinct whole-group supports across 42,384 support--missing-pair incidences in its root-hole sector. The maximum there at the accepted target was \(-1{,}968{,}374{,}370{,}320\). The tighter competitor was the envelope without root-hole credit.

\subsection{Trying to turn the surviving vertex into a construction}

The final priced maximizer illustrates why the finite dual and a constructive lower bound remain distinct. After restoring the preferred leaves required by its nonzero support, its odd coset contains the six residues \(1,9,17,21,27,33\). Five receive \(C\); the preferred residue \(27\) receives zero. The actual mean in that coset is \(5C/32=23921/64000\), close to \(3/8\). Its even-site mean vanishes. Before adding any price, the geometric quotient at this vertex is about \(10.20312816\%\). Its odd coset cannot supply the required even singleton. The priced score subtracts from the objective precisely because the preferred mass is missing and leaf \(9\) is full. Retaining only the raw objective would hide the failure.

A separate constructive control started from the active vertices for five nearby prices, restored their smallest preferred-child closures, and tried changing one retained class at a time. Across 86 distinct resulting supports, a real-valued solver proposed rates, which we rounded and checked against the exact total and the strict standard tile margins using rational logarithm enclosures. The best profile in this search has the same 28-leaf closure. Its two exceptional rates are \(0.70626403\) at residue \(27\) and \(1.09955197\) at residue \(46\); its other 26 occupied leaves all have rate \(2.392084\). The odd coset now has a genuine preferred mass. On this one word the ratio is
\[
 \frac{433565225}{4284671539}=10.1189839420\%.
\]
All 21 distinct pairs of cell means pass the strict rational checks. This attained value is a diagnostic for the adverse word. The separate all-label construction supplies the stronger lower bound quoted at the beginning of this section. A small repair of the closed vertex restores an admissible geometry, but this intermediate interval alone does not tell us whether a different support beats that repaired profile or whether more occupancy constraints can exclude the competing closed vertices. The final census in Section~\ref{sec:cycle64-endgame} separates these explanations by identifying every nominal support capable of exceeding the constructive guarantee.
